\VerdierPhysicalPageAnchor{0068}%
\VerdierSemanticAnchor{ch1-2-2-1-continuation}%
groups and that $\varepsilon(e_j,i)=-1$ if $j<i$,
$\varepsilon(e_j,i)=1$ if $i\leq j$. The map
$(\sigma,i)\longmapsto\varepsilon(\sigma,i)$ is moreover uniquely
determined by these properties. One immediately deduces the relation

\VerdierSemanticAnchor{ch1-2-2-1-eq2}
\begin{equation}
  \varepsilon(\sigma,i)\varepsilon(\sigma+e_i,j)
  =-\varepsilon(\sigma,j)\varepsilon(\sigma+e_j,i)
  \tag{2.2.1.2}
\end{equation}
valid for every $\sigma\in\mathbb Z[n]$ and every pair
$(i,j)\in[n]\times[n]$, $i\neq j$.

\VerdierSemanticAnchor{ch1-2-2-2}
\noindent\textbf{2.2.2.}\quad
Let $Y$ be an $n$-fold complex in $C$. Put

\VerdierSemanticAnchor{ch1-2-2-2-eq1}
\begin{equation}
  \delta_Y^{\sigma,i}=\varepsilon(\sigma,i)d_Y^{\sigma,i}.
  \tag{2.2.2.1}
\end{equation}

\VerdierSemanticAnchor{ch1-2-2-2-eq2}
It then follows from (2.2.1.2) that relations (2.1.2.1) are equivalent to
the relations
\begin{equation}
  \begin{aligned}
    \delta_Y^{\sigma+e_i,i}\delta_Y^{\sigma,i}
      &=0,
      &&\forall\sigma\in\mathbb Z[n],\quad\forall i\in[n],\\
    \delta_Y^{\sigma+e_j,i}\delta_Y^{\sigma,j}
      +\delta_Y^{\sigma+e_i,j}\delta_Y^{\sigma,i}
      &=0,
      &&\forall\sigma\in\mathbb Z[n],\quad\forall i,j\in[n].
  \end{aligned}
  \tag{2.2.2.2}
\end{equation}

\VerdierSemanticAnchor{ch1-2-2-2-pushforward}
Now let $\varphi\colon[n]\longrightarrow[m]$ be a map and suppose that
$Y$ is $\varphi$-summable. We then define an $m$-fold complex
$\int_\varphi Y$ by the formulas

\VerdierSemanticAnchor{ch1-2-2-2-eq3}
\begin{equation}
  \left(\int_\varphi Y\right)^\tau
  =\bigoplus_{\varphi(\sigma)=\tau}Y^\sigma,
  \qquad \forall\tau\in\mathbb Z[m],
  \tag{2.2.2.3}
\end{equation}

\VerdierSemanticAnchor{ch1-2-2-2-eq4}
\begin{equation}
  \begin{aligned}
    \delta_{\int_\varphi Y}^{\tau,i}(y_\sigma)
      &=\sum_{\varphi(j)=i}\delta_Y^{\sigma,j}(y_\sigma),\\[-1mm]
    &\hspace{-18mm}
      \sigma\in\mathbb Z[n],\quad i\in[m],\quad
      y_\sigma\in Y^\sigma,\quad\tau=\varphi(\sigma).
  \end{aligned}
  \tag{2.2.2.4}
\end{equation}

\VerdierSemanticAnchor{ch1-2-2-2-additive-interpretation}
Formula (2.2.2.4) defines a morphism from
$\left(\int_\varphi Y\right)^\tau
=\bigoplus_{\varphi(\sigma)=\tau}Y^\sigma$ to
$\left(\int_\varphi Y\right)^{\tau+e_i}
=\bigoplus_{\varphi(\sigma)=\tau+e_i}Y^\sigma$ when $C$ is the category
of abelian groups. It also makes it possible to define a morphism from
$\left(\int_\varphi Y\right)^\tau$ to
$\left(\int_\varphi Y\right)^{\tau+e_i}$ when $C$ is an additive category,
by working in the category of contravariant functors on $C$ with values
in the category of abelian groups, i.e. by interpreting $y_\sigma$ as a
morphism from an arbitrary object of $C$ to $Y^\sigma$.

\VerdierSemanticAnchor{ch1-2-2-2-verification}
An immediate calculation shows that the
$\delta_{\int_\varphi Y}^{\tau,i}$ satisfy relations (2.2.2.2), whence,
on returning to the differentials, the objects defined by (2.2.2.3)
equipped with the morphisms defined by

\VerdierSemanticAnchor{ch1-2-2-2-eq5}
\begin{equation}
  \begin{aligned}
    d_{\int_\varphi Y}^{\tau,i}(y_\sigma)
      &=\varepsilon(\tau,i)
        \sum_{\varphi(j)=i}
          \varepsilon(\sigma,j)d_Y^{\sigma,j}(y_\sigma),\\[-1mm]
    &\hspace{-18mm}
      \tau\in\mathbb Z[m],\quad i\in[m],\quad
      \sigma\in\mathbb Z[n],\quad\varphi(\sigma)=\tau
  \end{aligned}
  \tag{2.2.2.5}
\end{equation}
